% ======================= 译者序 =======================

\clearpage
\chapter*{译\ 者\ 序}
\markboth{译者序}{译者序}
\addcontentsline{toc}{chapter}{译者序}

二〇二一年十二月，我在莫斯科见到了巴兰尼科夫教授。那是一次原定半小时的会面，谈的是《战争的雷霆》中译本的修订，地点在列宁山他那间堆满微缩胶片的办公室。谈完正事之后，他从最底层铁皮柜里抽出一只牛皮纸袋，倒出三十七张黑白照片，摊在桌面上。

我看了二十分钟，什么也没看出来。照片是黑白的，物体是长条形的，表面刻着字，字很小，笔画匀细。我认出几个，其余不认识。我以为是某种篆书的变体，说了出来。他没有否认，也没有承认，只把其中一张卡片推到我面前。卡片是俄文的，打印在一张发黄的硬纸片上，编号一栏写着：\textsc{фв}-1963-11，下面一行汉字，是译文：「第 4 号居住环」。我问他这是什么地方。他说：月球正面，静海西缘，一处编号 \textsc{mare-tranq-07} 的基地。那是我第一次听说广寒宫不是一座宫殿。

这三十七张照片拍于一九六三年。此后六十年间，这一叠卡片在莫斯科那间档案室里没有被第二个人调阅过。巴兰尼科夫从二十五岁那年开始整理它们，到我见他那一年，整整五十八年。他自己称这一卷为「我一辈子做得最慢的事」。

\section*{一、\ 关于这一卷的位置}

按照丛书的总目录，《月之都》排在第一卷，出版时间却在第七卷《战争的雷霆》之后。这并非编辑上的失误，而是巴兰尼科夫本人的安排。他一直认为，先把正史里那些说不通的地方摆出来，再把这批没有落款的残片端上桌面，读者才不会把它们当作伪造的古董看。

读者若是先读过《战争的雷霆》，会在本书里见到几个熟面孔：静海七号基地、第 117 全宗第 2187 号卷宗所录退役清单的十九项处置、六十年前苏联方面「来源不明，暂不归类」的编目结论，以及那一群自月球分六批迁回地面、自此不复通音书的人。需要说明的是，两本书的关系不是「互证」，也不是谁给谁作注：本书直接据莫斯科档案馆第 117 全宗、VFS 基金会编目及巴兰尼科夫此前两部俄文专著行文，材料同一而体例各异。《战争的雷霆》是一部概论，用各路档案拼出一座城的轮廓；本书则只做两件事：把这三十七片能读出来的字按顺序摆好，再把它们能说明与不能说明的东西一一分清。

\section*{二、\ 关于「辑佚」二字}

中文书名里的「辑佚与考释」五个字，是张建国和我共同商定的，俄文原题直译应作《月之都：广寒宫残片三十七种整理与研究》。原题过于朴素，放到中文学术书目里容易被误认为考古报告。用「辑佚」，取的是中国古典文献学的意思：散佚之文，重新集录；但需要声明，本书辑的不是纸上散佚的书，而是三十七件留在一间仓库里的实物。

翻译上最费斟酌的是上编的三十七条释文。巴兰尼科夫的俄文版是将小篆转写为现代楷书汉字，再附俄文直译。中文版若再从俄文转译回来，三十七条就成了二次翻译，韵味全失；若直接采用他转写出来的楷书，文字又过于简古，普通读者未必顺畅。我们最后定的办法是：释文一律照原样保留，一个字也不改，连补字符号也照原样；另外在书末加一篇附录，把三十七片逐条译成白话，读者读不通时翻到后面查。这不是一个好办法，它让读者读一段要翻两回书。我们试过别的方案，比如把白话夹在释文里，或者把释文整体改写成现代汉语，都不行：前者破坏了原件的样式，后者等于替巴兰尼科夫重写了一遍。三十七片的译文放在附录四，由张建国初译、我复校，出错算我们的。

\section*{三、\ 关于那枚圆片}

本书简三十六所记「一物，圆，径不及二寸……击之不鸣」，即存藏于基金会、编号 \textsc{фв}-1963-46 的那枚金属圆片。这枚圆片我们在《战争的雷霆》的译者序里已经提过一次，当时是因为它的铼锇同位素比值；这一次重提，是因为二〇二三年我们请人对它作了一次表面微区分析，发现表面的刻痕是一圈极细的同心环，不是装饰，环距不等，最密处在每毫米一百二十条左右。这个密度超出了任何已知的机械刻写方式。负责检测的实验室在报告末页加了一句手写附注：「建议送光学机构复测，本所不负责解释此种加工方法。」我们没有再取样，这一句附在这里，权作存录。

\section*{四、\ 体例说明}

本书体例较《战争的雷霆》繁复，需要预先交代三事。一、上编释文一律以楷书排印，小篆原貌须待第三版影印本附出；释文中的标点为辑校者所加，原件无句读。二、补字作方框（□），缺字不可知者亦作方框，两者在正文里看不出区别，须看校记；这是原书的做法，我们没有改。三、书中凡引及「第 117 全宗」「\textsc{фв}-」号卷宗与俄方先期研究成果，分别据莫斯科档案馆原始编目、VFS 基金会档册及巴兰尼科夫《月面设施复原图集》（2016）、《Лунная прародина》（2018）。

\section*{五、\ 一件没有解决的事}

最后说一件事，说出来对本书不利。翻译过程中，我们核对了全书所有的数字，凡能互相校验的都校过，唯有一处始终对不上：简十九记事故二万二千，简二十记实存十一万九千，而按莫斯科档案馆第 117 全宗第 2187 号卷宗的先期复原，城池完工时的人数是十五万。两数之间差了八千。巴兰尼科夫在书里承认了这个差额，并给了三种可能的解释，最后坦承倾向前者但没有证据。我们照译，没有替他补。这一处之所以要在这里点出来，是因为它是全书唯一一处作者自己承认不知道的地方：其余的判断他都很确定，只有这一处他写明没有证据，我们认为这一段应当被中文读者看到。

\begin{flushright}
{\small 李文亚}\quad{\footnotesize\color{pencil} 二〇二四年七月，于北京}
\end{flushright}


% ======================= 中文版序言 =======================

\clearpage
\chapter*{中文版序言}
\markboth{中文版序言}{中文版序言}
\addcontentsline{toc}{chapter}{中文版序言}

这本书是为一群我从未见过的人写的。三十七片残简上的字，是汉字；读它们的人，在中国。而把它们从月面带回地球、编目、封存、六十年无人过问的，是另一群人。我做的事，只是把两头连起来。连起来之后，我发现一个尴尬的事实：这些字写得很坏，坏在它们没有文采，没有修辞，没有一句多余的话，记账的只记账，登记的只登记，连最后一天的场面也只用了二十七个字。若把它们当作文章来读，读三页就会困。但我建议各位读下去，理由只有一个：这是一个月球上的城在最后两万年里，唯一愿意留给自己的东西。一个有十一万九千人的城，两万年，只留下三十七片，平均五百四十年一片；这个比例本身，比任何一片的内容都更能说明那座城是什么样子。

\section*{一、\ 关于我的读法}

中国读者大约会问：一个俄国人，凭什么读小篆。这个问题我答过很多次，答案一直不太好听：凭笨功夫。我从一九六三年开始抄这些字，抄了六十年，抄错的地方不少，这我知道；但凡是能互相校验的地方，我都校过三遍。至于方法，我在导论里写了五节，此处不重复，只补一句：这本书里凡是我说有把握的地方，都是因为那一片的字迹清楚、上下文能对上、而且能与其他几片互证；凡是这三条缺一条的，我都写了「残简无一句」。这句话全书出现七次。这话写在纸上，是记录，不是场面话。

\section*{二、\ 关于这本书的命运}

最后说一件与学术无关的事。这一卷在莫斯科印了两次，第一次是一九八〇年代的内部油印本，印数不足百册；第二次是现在这个排印本。两次之间隔了四十年。四十年里，这批材料没有丢，原因不好听：它被放在库房最里面，几次清退都没轮到。我把这件事写进结语，不是发牢骚；我想说的是，一本书能不能留下来，与它写得好不好关系不大。这一点，中国的读者大概比我更清楚。

\begin{flushright}
{\small А.\,В.\,巴兰尼科夫}\quad{\footnotesize\color{pencil} 二〇二四年八月，于莫斯科}
\end{flushright}
